% ADDED (2026-11-02): thread 1 of the conclusions, inertia vs anchor.
\paragraph{Inertia and anchor.}
An agent can resist the opinions of the others in two ways, one for each model.
In the FDG model, a self-weight $a_{ii}$ is inertia:
% REVISED (2026-11-05, review, R19). Was:
% A self-weight $a_{ii}$ is inertia:
an agent whose neighbours all hold the opinion $1$ updates as
$x_i(k+1) = a_{ii}\, x_i(k) + (1 - a_{ii})$ and reaches $1$, only more slowly as $a_{ii}$
grows. In the FJ model, an attachment $1 - \lambda_i$ is an anchor: the same agent updates as
% REVISED (2026-11-05, review, R19). Was:
% grows. An attachment $1 - \lambda_i$ is an anchor: the same agent updates as
$x_i(k+1) = \lambda_i + (1 - \lambda_i)\, u_i$ and stops at once between $1$ and its
prejudice, never reaching the others. This holds for a single agent on any network; the
effects on the group are shown on the two networks of Figure~\ref{fig:networks}, both
strongly connected and periodic without self-loops.
\begin{itemize}
  \item \emph{Convergence.} In the FDG model the self-weight of agent~$1$ is needed to
        make the network aperiodic (Section~\ref{subsec:fdg-stub}); in the FJ model a
        single prejudiced agent that influences everybody makes all agents P-dependent, and
% REVISED (2026-11-04, editorial, S1). Was:
%         single prejudiced agent reached by everybody makes all agents P-dependent, and
        the model converges although the network is periodic
        (Corollary~\ref{cor:fj-stable}).
  \item \emph{Influence.} In the FDG model, the social power of agent~$1$ grows gradually with its
% REVISED (2026-11-05, review, R19). Was:
%   \item \emph{Influence.} The social power of agent~$1$ grows gradually with its
        self-weight, from $1/6$ to $1$ (Figure~\ref{fig:fdg-power}), and for $t < 1$
        the consensus still depends on the initial opinions of all agents. In the FJ
        model, a single anchor, even a weak one,
% REVISED (2026-11-05, review, R19). Was:
% the consensus still depends on the initial opinions of all agents. A single
%         anchor, even a weak one,
 takes the whole influence centrality,
% REVISED (2026-11-04, editorial, O5). Was:
% A single
%         anchor, however loose, takes the whole influence centrality,
 $c = e_1$, and the
        group ends on its prejudice (Figure~\ref{fig:fj-one-prejudiced}).
  \item \emph{Many agents.} With the same weight $0.2$ for every agent
        (Figure~\ref{fig:fj-all-prejudiced}), the FDG model reaches a consensus, as on
        every strongly connected network with a self-loop
        (Corollary~\ref{cor:primitive-selfloop}), while the FJ model keeps a disagreement
        as soon as the prejudices differ, on any network.
\end{itemize}
In social terms, inertia buys time and a gradual influence, but in the end the agent
follows the group; an anchor buys permanence: the agent never gives in completely, a
single anchor takes all the influence, and different anchors keep the group divided.

% ADDED (2026-11-02): thread 2 of the conclusions, memory in the state vs external input.
\paragraph{State and input.}
The same opinion can be a state or an input, and this decides the stability. An FDG
model with stubborn agents and its FJ rewriting~\cite{proskurnikov2017tutorial} have the
same trajectories, but the FDG model is never asymptotically stable, while the FJ model
is asymptotically stable whenever the stubborn agents influence all the others
(Corollaries~\ref{cor:fdg-stubborn} and~\ref{cor:fj-stable}): the opinion
% REVISED (2026-11-04, A4). Was:
% same trajectories, but the FDG model is never asymptotically stable, while the FJ model
% is (Corollary~\ref{cor:fj-stable}): the opinion
 of a stubborn agent is a state in the
first and the input $u_i = x_i(0)$ in the second. The spectra show it in one step: the
stubborn row is $e_i^\top$ in $A$ and zero in $\Lambda W$, so the eigenvalue $1$ of the
FDG model becomes $0$ in the FJ model. Conversely, the augmented
model~\eqref{eq:fj-augmented} turns every prejudice back into the state of a virtual
agent, adds one eigenvalue $1$ for each of them, and loses the asymptotic stability.
\begin{itemize}
  \item \emph{What is perturbed.} Moving $x_1(0)$ (Figure~\ref{fig:fj-stub-perturbation})
        moves a state in the FDG model, which keeps it forever, and only the initial
        opinion of a real agent in the FJ model, which forgets it and returns to the same
        limit $Vu$. Moving $u_1$ would change the FJ limit as well.
  \item \emph{Memory without stubbornness.} A closed group of P-independent agents is
        a memory inside the real agents (Figure~\ref{fig:fj-auxiliary}): it keeps its
        initial opinions as a virtual agent keeps its prejudice, so the FJ model converges
        without being asymptotically stable.
\end{itemize}
In social terms, what a group remembers is what the model counts as its state: the
opinion of a stubborn agent in the FDG model, the initial opinions of a closed group of
P-independent agents in the FJ model, but never the initial opinions of the agents that a
prejudice reaches.
% REVISED (2026-11-04, editorial, S3). Was:
% In social terms, in the FDG model stubbornness is a memory inside the group, which is
% never forgotten; in the FJ model it is an external influence, and the group forgets its
% initial opinions: what the group remembers is what the model counts as its state.

% ADDED (2026-11-02): thread 3 of the conclusions, social power vs influence centrality.
\paragraph{Social power and influence centrality.}
Both measures say how much the initial opinion of each agent weighs in the final
opinions, under the different assumptions listed in Table~\ref{tab:fdg-fj}; the influence
centrality also averages over the agents, since the rows of $V$ differ.
% REVISED (2026-11-04, editorial, S4). Was:
% Both measures say how much the initial opinion of each agent weighs in the final
% opinions, but under different assumptions. The social power $w$~\eqref{eq:social-power}
% requires a consensus, $A^\infty = \mathbf{1}_n w^\top$. The influence centrality
% $c = n^{-1} V^\top \mathbf{1}_n$~\eqref{eq:influence} requires $\rho(\Lambda W) < 1$ and
% $u = x(0)$, which gives the initial opinions a weight in $x(\infty) = V x(0)$, and it
% averages over the agents, since the rows of $V$ differ.
\begin{itemize}
  \item \emph{A single anchor.} With agent~$1$ the only prejudiced agent, $c = e_1$ for
        every $\lambda_1 < 1$, while in the FDG model $w_1$ grows gradually with the
        self-weight. The augmented model explains the value: without the virtual agents
        that nobody listens to, it has a single sink $1'$, reaches a consensus and has
        social power $e_{1'}$, so with a single prejudice the two measures coincide.
  \item \emph{Many anchors.} With $\Lambda = \alpha I$ every agent is anchored, the
        augmented model has $n$ sinks that are listened to and no social power, while
        $c_\alpha$ exists for every $W$. By Lemma~\ref{lem:centrality-limit},
        $c_\alpha \to w$ as $\alpha \to 1$ when $W$ is fully regular
        (Figure~\ref{fig:fj-lemma24}): the influence centrality moves from the uniform
        value $\mathbf{1}_n/n$, when every agent keeps its own opinion, to the social
        power, when the anchors vanish.
\end{itemize}
In social terms, the social power tells who leads a group that agrees, the influence
centrality who weighs most in a group that keeps disagreeing; the second becomes the
first when the anchors are released.

% ADDED (2026-11-03): part C of the conclusions, what is left out.
\paragraph{What is left out.}
In Section~\ref{subsec:fj-stub} the prejudices are the initial opinions, $u = x(0)$, as in
the original theory of Friedkin and Johnsen, apart from the perturbations of $x(0)$ in
Figures~\ref{fig:fj-stub-perturbation} and~\ref{fig:fj-auxiliary}(d),
% REVISED (2026-11-04, B5: Figures 7 and 12(d) have u different from x(0)). Was:
% In Section~\ref{subsec:fj-stub} the prejudices are the initial opinions, $u = x(0)$, as in
% the original theory of Friedkin and Johnsen,
 while~\cite{proskurnikov2017tutorial}
treats $u$ and $x(0)$ as independent: the examples could be repeated with
$u \neq x(0)$, for instance by perturbing all the prejudices slightly.
\cite{proskurnikov2017tutorial} also reads the final opinions of the FJ model as a Nash
equilibrium of a game in which each agent balances its distance from the others and from
its prejudice, and discusses the continuous-time models of Abelson and Taylor; both are
beyond the scope of this report.
